% Source: https://artofproblemsolving.com/community/c3h1494409_drawing_altitudes_from_all_sides_of_a_triangle
% Stable identity: aops-topic-1494409-q01
\textbf{(a)} The foot $D$ is not $C$, since angle $C$ is not a right angle. If $D$ were on ray $CB$, then triangle $ACD$ would have the right angle at $D$ and also the angle $\angle ACD=\angle ACB>90^\circ$, which is impossible (the angles of a triangle add to $180^\circ$). So $D$ is on the ray opposite to $CB$, that is, on line $BC$ beyond $C$. In the same way $E$ is on line $AC$ beyond $C$.

\textbf{(b)} Lines $AD$ and $BE$ are not parallel (they are perpendicular to the non-parallel lines $BC$ and $AC$), so they meet at a point $H$. Write points as vectors. Since $AH\perp BC$ and $BH\perp CA$,
\[
(H-A)\cdot(B-C)=0,\qquad (H-B)\cdot(C-A)=0 .
\]
Adding, the terms $H\cdot B$, $H\cdot C$ and $A\cdot B$ cancel in pairs and we get
\[
H\cdot B-H\cdot A+A\cdot C-B\cdot C=(H-C)\cdot(B-A)=0 .
\]
So $CH\perp AB$, that is, $H$ lies on line $CF$ as well. The three altitude lines are concurrent.

By (a), line $AD$ meets line $BC$ at $D$, which is not on side $BC$. A line through the vertex $A$ that passes through any other point of the triangle must cross the opposite side $BC$; so line $AD$ meets the triangle only at $A$. Also $H\ne A$: otherwise $A$ would lie on the altitude line from $B$, giving $BA\perp AC$, a right angle at $A$, impossible since the angle at $C$ is already obtuse. Hence $H$ lies outside the triangle (beyond $C$, as in the figure).

\begin{asy}
usepackage("amsmath");
size(6.5cm);
defaultpen(fontsize(10pt));
pair ft(pair P,pair U,pair V){return U+dot(P-U,V-U)/dot(V-U,V-U)*(V-U);}
path ra(pair P,pair u,pair v,real s=0.2){return P+s*unit(u)--P+s*unit(u)+s*unit(v)--P+s*unit(v);}
pair A=(0,0),B=(6,0),C=(4.2,1.8);
pair D=ft(A,B,C),Ee=ft(B,A,C),F=ft(C,A,B);
pair H=extension(A,D,B,Ee);
draw(C--D,gray+dashed); draw(C--Ee,gray+dashed);
draw(A--H,blue+dashed); draw(B--H,blue+dashed); draw(F--H,blue+dashed);
draw(A--B--C--cycle,linewidth(0.9));
draw(ra(D,C-D,A-D)); draw(ra(Ee,C-Ee,B-Ee)); draw(ra(F,B-F,H-F));
dot(A);dot(B);dot(C);dot(D);dot(Ee);dot(F);dot(H,red);
label("$A$",A,SW); label("$B$",B,SE); label("$C$",C+0.5*dir(335));
label("$D$",D,W); label("$E$",Ee,E); label("$F$",F,S); label("$H$",H,N,red);
label("\mbox{}",(min(currentpicture,user=true)+max(currentpicture,user=true))/2,black); // keeps the text after the figure black when LaTeX typesets the labels inline
\end{asy}

\textbf{(c)} In triangle $ABH$ the line through $A$ perpendicular to $BH$ is line $AC$ (because $BH\perp AC$), the line through $B$ perpendicular to $AH$ is line $BC$ (because $AH\perp BC$), and the line through $H$ perpendicular to $AB$ is line $HC$. All three pass through $C$, so $C$ is the orthocentre of triangle $ABH$. $\blacksquare$

So for an obtuse triangle the altitudes themselves do not meet; it is the lines containing them, extended beyond the obtuse vertex, that meet at $H$.
